\documentclass[11pt]{article}
\newcommand{\DocShort}{Examen: hidrostática y propiedades seleccionadas}
\input{exam_preamble.tex}

\begin{document}

\doccover{Examen de hidrostática y propiedades seleccionadas de los fluidos}{2026-07}
\introbox{Evaluar la capacidad del estudiante para resolver problemas de hidrostática y de propiedades básicas de los fluidos mediante el uso correcto de principios físicos, equilibrio y relaciones constitutivas simples.}{Se espera una solución clara, ordenada y completamente trazable. En cada ejercicio el estudiante deberá realizar primero un dibujo o esquema del sistema, identificar datos y variables, plantear las ecuaciones aplicables y desarrollar la solución en unidades SI.}

\section{Objetivo del examen}
Resolver seis problemas representativos de presión hidrostática, manometría, empuje sobre superficies, flotación, estabilidad de presas, tensión superficial y viscosidad.

\section{Instrucciones generales}
\begin{itemize}
\item Resuelva en unidades SI.
\item Presente claramente las ecuaciones utilizadas y las sustituciones numéricas.
\item Puede asumir $g = \SI{9.81}{\meter\per\second\squared}$.
\item Cuando sea necesario, use para el agua a \SI{20}{\celsius}: $\rho = \SI{998}{\kilogram\per\meter\cubed}$ y $\gamma = \SI{9.79}{\kilo\newton\per\meter\cubed}$.
\item Valor total: \SI{100}{} puntos.
\end{itemize}

\begin{center}
\begin{tabular}{lc}
\toprule
\textbf{Problema} & \textbf{Puntuación sugerida} \\
\midrule
1. Presión hidrostática y manometría & 15 puntos \\
2. Empuje hidrostático sobre superficie plana & 20 puntos \\
3. Flotación y principio de Arquímedes & 15 puntos \\
4. Estabilidad de una presa de gravedad & 25 puntos \\
5. Tensión superficial & 12.5 puntos \\
6. Viscosidad & 12.5 puntos \\
\bottomrule
\end{tabular}
\end{center}

\section{Datos y constantes de uso común}
\begin{itemize}
\item Presión atmosférica de referencia, cuando aplique: $p_{\text{atm}} = \SI{101.3}{\kilo\pascal}$.
\item Para el mercurio, cuando aplique: $\rho_{\text{Hg}} = \SI{13600}{\kilogram\per\meter\cubed}$.
\item Para el concreto de la presa, cuando aplique: $\gamma_c = \SI{23.6}{\kilo\newton\per\meter\cubed}$.
\item Para la solución jabonosa del problema de tensión superficial: $\sigma = \SI{0.025}{\newton\per\meter}$.
\end{itemize}

\section{Problemas}

\subsection*{Problema 1. Presión hidrostática y manometría}
\textbf{Antes de iniciar, realice un esquema claro del sistema, indicando niveles, fluidos, puntos de referencia y datos geométricos.}

Un tanque cerrado contiene agua y aire comprimido. La presión del aire sobre la superficie libre del agua se mide mediante un manómetro en U conectado al espacio de aire del tanque. El otro extremo del manómetro está abierto a la atmósfera y contiene mercurio.

La diferencia entre los niveles de mercurio en el manómetro es de \SI{18}{\centi\meter}, siendo el nivel del lado conectado al tanque más bajo que el lado abierto. La profundidad del agua dentro del tanque, medida desde la superficie libre hasta el punto $A$ en el fondo, es de \SI{2.40}{\meter}.

Determine:
\begin{enumerate}
\item La presión manométrica del aire dentro del tanque.
\item La presión absoluta del aire dentro del tanque.
\item La presión manométrica en el punto $A$.
\item La presión absoluta en el punto $A$.
\item Explique brevemente por qué la presión absoluta y la presión manométrica difieren, y cuál de ellas resulta más útil para describir el estado del agua en el fondo del tanque.
\end{enumerate}

\subsection*{Problema 2. Empuje hidrostático sobre superficie plana}
\textbf{Antes de iniciar, realice un esquema claro del sistema, indicando dimensiones, nivel del agua, articulación, apoyo y línea de acción de la resultante.}

Una placa rectangular vertical de \SI{5.0}{\meter} de ancho y \SI{4.0}{\meter} de altura cierra el extremo de un canal de agua dulce. La lámina de agua alcanza una profundidad de \SI{3.2}{\meter} medida desde la superficie libre hasta el fondo del canal. La placa está articulada a lo largo de su borde superior en el punto $A$ y se apoya contra un tope en el punto $B$, ubicado en su borde inferior.

Determine:
\begin{enumerate}
\item La magnitud de la fuerza hidrostática resultante ejercida por el agua sobre la placa.
\item La posición del centro de presión, medida verticalmente desde la superficie libre.
\item La fuerza horizontal que debe ejercer el tope en $B$ para mantener la placa en equilibrio.
\end{enumerate}

Desprecie el peso de la placa y la fricción en la articulación.

\subsection*{Problema 3. Flotación y principio de Arquímedes}
\textbf{Antes de iniciar, realice un esquema claro del sistema e indique las fuerzas relevantes para el análisis.}

Se desea determinar el volumen y la densidad promedio de una pieza sólida irregular empleando una balanza de resorte. La pieza pesa \SI{6.80}{\kilo\newton} en el aire y \SI{4.25}{\kilo\newton} cuando está completamente sumergida en agua a \SI{20}{\celsius}.

Determine:
\begin{enumerate}
\item El volumen de la pieza.
\item La densidad promedio de la pieza.
\item La gravedad específica del material.
\item Si la pieza pudiera colocarse libremente en agua, indique si flotaría o se hundiría. Justifique su respuesta con base en los resultados obtenidos.
\end{enumerate}

Exprese claramente sus suposiciones.

\subsection*{Problema 4. Estabilidad de una presa de gravedad}
\textbf{Antes de iniciar, realice un esquema claro de la presa, indicando geometría, fuerzas, centro de presión y base de análisis.}

Una presa de gravedad de concreto, de sección rectangular, tiene una altura de \SI{18}{\meter} y una base de \SI{12}{\meter}. La cara aguas arriba es vertical y el embalse está lleno hasta la coronación de la presa. Considere un ancho unitario de \SI{1.0}{\meter} perpendicular al plano del análisis.

Para el agua, use $\gamma_w = \SI{9.79}{\kilo\newton\per\meter\cubed}$.

Determine:
\begin{enumerate}
\item La magnitud de la fuerza hidrostática ejercida por el agua sobre la cara vertical de la presa.
\item La ubicación del centro de presión sobre la cara mojada, medida desde la superficie libre.
\item El peso propio de la presa por metro de ancho.
\item La magnitud de la resultante entre la fuerza hidrostática y el peso de la presa.
\item El punto donde la resultante corta la base de la presa, medido desde el talón aguas arriba.
\item Si la resultante pasa dentro del tercio central de la base.
\item Con base en el resultado anterior, concluya si existe riesgo de tracción en la base y si la presa es estable frente al volcamiento en esta condición.
\end{enumerate}

Desprecie subpresión, sismo, oleaje, sedimentos y cualquier otra fuerza adicional.

\subsection*{Problema 5. Tensión superficial}
\textbf{Antes de iniciar, realice un esquema claro del sistema y represente la dirección de la fuerza de tensión superficial.}

Un aro circular delgado de alambre, de \SI{8.0}{\centi\meter} de diámetro, se sumerge en una solución jabonosa y luego se extrae lentamente, de modo que se forma una película líquida sobre el aro.

Determine:
\begin{enumerate}
\item La fuerza mínima que debe aplicarse verticalmente para sostener la película justo antes de desprenderse del aro.
\item El trabajo necesario para elevar uniformemente la película una distancia de \SI{6}{\milli\meter}, suponiendo que la fuerza permanece constante durante ese desplazamiento.
\item Explique por qué en este tipo de problema debe considerarse que la película tiene dos superficies libres.
\end{enumerate}

\subsection*{Problema 6. Viscosidad}
\textbf{Antes de iniciar, realice un esquema claro del sistema, indicando área de contacto, espesor de película, dirección del movimiento y perfil de velocidad supuesto.}

Un bloque rectangular de \SI{0.50}{\meter} de largo, \SI{0.30}{\meter} de ancho y peso despreciable se mueve a velocidad constante sobre una superficie horizontal, separado de ella por una película de aceite de espesor uniforme \SI{0.80}{\milli\meter}. La velocidad del bloque es de \SI{0.60}{\meter\per\second}.

Si la fuerza horizontal necesaria para mantener el movimiento es de \SI{42}{\newton}, determine:
\begin{enumerate}
\item El esfuerzo cortante promedio en la película de aceite.
\item La viscosidad dinámica del aceite.
\item La viscosidad cinemática del aceite, si su densidad es $\rho = \SI{870}{\kilogram\per\meter\cubed}$.
\item Si el espesor de la película se reduce a la mitad y todo lo demás permanece constante, indique cómo cambia la fuerza requerida para mantener la misma velocidad.
\end{enumerate}

Suponga distribución lineal de velocidades en la película.

\section{Rúbrica general}
\begin{infobox}
Su calificación en cada problema dependerá de la calidad integral de su solución. En términos generales, se valorará lo siguiente:
\begin{itemize}
\item que inicie cada ejercicio con un dibujo o esquema claro del sistema;
\item que identifique correctamente los datos, incógnitas y suposiciones;
\item que plantee las ecuaciones físicas apropiadas para el problema;
\item que desarrolle los cálculos con orden, unidades consistentes y sustituciones claras;
\item que obtenga resultados numéricos razonables e interprete su significado cuando se le solicite.
\end{itemize}
\end{infobox}

\begin{center}
\begin{tabularx}{\textwidth}{>{\raggedright\arraybackslash}p{0.52\textwidth} >{\raggedleft\arraybackslash}p{0.16\textwidth} >{\raggedright\arraybackslash}X}
\toprule
\textbf{Criterio} & \textbf{Peso guía} & \textbf{Qué debe mostrar} \\
\midrule
Esquema y organización inicial & 15\% & Dibujo claro, definición de variables, datos completos y lectura correcta del enunciado. \\
Planteamiento físico y ecuaciones & 35\% & Uso correcto de principios, equilibrio, relaciones constitutivas y ecuaciones aplicables. \\
Desarrollo matemático y unidades & 30\% & Procedimiento ordenado, sustituciones claras y manejo consistente de unidades SI. \\
Resultado e interpretación & 20\% & Resultado final correcto o razonable, con interpretación cuando el ejercicio la requiera. \\
\bottomrule
\end{tabularx}
\end{center}

\end{document}
